\originalsection{18.315：作业2（原卷 Fall 2005）}
本组收录原题2、3、4, 共四个任务单元. 这里的平面三角剖分为简单平面图,
包括外面在内的每个面都是三角形; 从图中减去一个面, 指删除其三条边, 保留顶点.

\begin{exercise}\label{mit:18315-hw2-q2}
\textbf{18.315, 原卷 Fall 2005, 作业2, 题2.}
爪为 $K_{1,3}$. 爪覆盖是把图的全部边分成三条边一组,
每组的三条边共享一个中心顶点. 证明每个平面三角剖分 $T$
都存在一个面 $F$, 使 $T-E(F)$ 有爪覆盖.
\end{exercise}
\sourceof{官方 HW2 第1页题2; 编者独立解答}
\begin{solution}
实际上任意一个面都可以. 选 $F$ 为外面, 它的三个顶点记作 $O$,
其余顶点集为 $I$, 令 $G=T-E(F)$.
建立一个二部图: 左侧每条 $G$ 的边是一个顶点;
右侧每个 $v\in I$ 有三个槽位. 边顶点与其内端点的三个槽位相邻.
我们要把每条边分给一个内端点, 使每个内顶点恰收到三条边.

核 Hall 条件. 对任意一组边 $B$, 设 $S\subseteq I$ 为它们的全部内端点.
这些边全在 $G[S\cup O]$ 中. 将外三角形的三条边加回,
得到一个有 $|S|+3$ 个顶点的简单平面图, 因而边数至多
$3(|S|+3)-6=3|S|+3$. 删除那三条边后,
\[
|B|\le |E(G[S\cup O])|\le3|S|=|N(B)|.
\]
Hall 定理给出覆盖全部左侧边顶点的匹配.
三角剖分有 $3|V(T)|-6$ 条边, 故
$|E(G)|=3|V(T)|-9=3|I|$, 左右两侧等大,
每个内顶点的三个槽位都用完. 把分给同一顶点的三条边组成一个爪,
就得到所需覆盖. 当 $T$ 只有外三角形时, $G$ 无边, 空覆盖满足结论.
\end{solution}

\begin{exercise}\label{mit:18315-hw2-q3}
\textbf{18.315, 原卷 Fall 2005, 作业2, 题3.}
设 $G=T-E(F)$, $F$ 是平面三角剖分的外面. 对一个爪覆盖,
把每个爪的三条边从中心向外定向. 翻转一个有向圈的全部边称为一次翻转.
原题要求证明: (a) 任意两个爪覆盖可由一列翻转相连;
(b) 即使只允许长度3的有向圈翻转, 结论仍成立.
\end{exercise}
\sourceof{官方 HW2 第1页题3; 原题未限制各顶点出度, 下述完整反例由编者独立构造并核验}
\begin{solution}
按题2给出的\textbf{任意爪覆盖}定义, 两个断言均不成立.
一次有向圈翻转在圈上每个顶点交换一条入边和一条出边,
所以每个顶点的出度都是不变量. 爪覆盖的定义却没有固定这个出度序列.

取二十面体图. 顶点为 $N=0,S=11$,
$a_i=1+i,b_i=6+i$（$i$ 按模5计算）. 六类边为
\[
Na_i,\quad Sb_i,\quad a_ia_{i+1},\quad b_ib_{i+1},\quad a_ib_i,\quad a_ib_{i-1}.
\]
两条五边形圈之间的三角形带加上上下两个锥盖给出球面三角剖分;
任选面 $Na_0a_1$ 为外面, 删除 $01,02,12$, 剩下27条边.
下表列出两个覆盖: 每个非空格中的三个数字是相应中心的三个叶顶点,
即该行代表三条由中心连出的边. 顶点0、1在两覆盖中都没有爪.
\begin{center}
\begin{tabular}{ccl}
\toprule
爪中心 & 覆盖 $A$ 的三个叶点 & 覆盖 $B$ 的三个叶点\\
\midrule
2 & 无 & $7,6,3$\\
3 & $4,2,0$ & $8,7,0$\\
4 & $5,8,0$ & $8,3,0$\\
5 & $10,1,0$ & $4,1,0$\\
6 & $10,2,1$ & $10,7,1$\\
7 & $6,3,2$ & 无\\
8 & $9,7,3$ & $11,9,7$\\
9 & $4,11,5$ & $4,10,5$\\
10 & $11,9,1$ & $11,5,1$\\
11 & $8,7,6$ & $6,9,7$\\
\bottomrule
\end{tabular}
\end{center}
可逐边检查: 每列的九组爪均无重复, 恰好覆盖全部27条边.
但顶点2在 $A$ 中出度0、在 $B$ 中出度3;
顶点7的情形相反. 因出度不变量, 两者连任意长度的圈翻转都不能相连,
更不用说只允许三角形翻转.

原题图示所暗示的常见版本另加外顶点出度0、内顶点出度3,
这是不同的条件. 对同一固定出度序列, (a)的圈翻转结论确实成立:
将两个定向中方向不同的边按第一个定向取出,
每点入度等于出度, 因两完整定向出度相同.
这个平衡有向图的边可逐次抽出有向圈, 翻转后相应的不同边消失,
有限步达到第二个定向. 此修订只解释缺失条件, 不替代原题的反例;
也不据此宣称已证明修订版本的三角形翻转结论.
\end{solution}

\begin{exercise}\label{mit:18315-hw2-q4}
\textbf{18.315, 原卷 Fall 2005, 作业2, 题4.}
在三角格点中取边长 $n$ 的等边三角形区域 $Q_n$, 把全部格点染成1、2、3三色,
不要求正常染色. 设 $f$ 为顺时针颜色循环 $1,2,3$ 的小三角形数,
减去顺时针颜色循环 $3,2,1$ 的小三角形数.
证明 $f$ 只取决于边界格点的颜色.
\end{exercise}
\sourceof{官方 HW2 第1--2页题4; 编者独立解答}
\begin{solution}
给有向格边赋权: 颜色变化 $1\to2,2\to3,3\to1$ 赋 $1/3$,
反向赋 $-1/3$, 同色赋0. 一个小三角形沿顺时针边界的权重和,
三色按正循环时为1、反循环时为 $-1$, 少于三色时为0:
双色三角形中的两条异色边恰互为反向, 权重抵消.
因此 $f$ 等于所有小三角形顺时针边界权重的总和.
每条内部边出现两次且方向相反, 全部抵消;
最终只剩整个 $Q_n$ 的顺时针边界权重和.
这一表达式只用边界颜色, 所以固定边界后改变内部任意格点颜色都不改变 $f$.
\end{solution}
